\newcommand{\UPSTItitreEnTete}{Chaîne de conversion -- Cellule robotique}
\newcommand{\UPSTItitre}{De la donnée NOC au bloc fonctionnel : comprendre la chaîne de conversion}
\newcommand{\UPSTItitreEnTetePages}{Chaîne de conversion NOC -- Code applicatif}
\newcommand{\UPSTIsousTitreEnTete}{Préparation -- HMI et robot}

% Trou de complétion inséré DANS un listing, au point exact où l'étudiant
% doit écrire (même principe que dans le DS1 S1) : escapeinside dédié à ce
% document, avec un délimiteur `(*@...@*)` qui ne peut pas être confondu
% avec un commentaire ST normal `(* ... *)`.
\lstset{escapeinside={(*@}{@*)}}
% Syntaxe : \trou[reponse]{largeur}
%   - version élève (\documentVersion{E}) : rectangle vide (réponse ignorée) ;
%   - version prof  (\documentVersion{P}) : réponse affichée en rouge dans
%     le rectangle (via \UPSTIcorrection), qui s'élargit si nécessaire.
% Dans un listing, la réponse est en mode LaTeX : échapper les « _ » (\_).
\newcommand{\trou}[2][]{%
    \tikz[baseline=(trou.base)]{%
        \node[draw=gray!70, dashed, rounded corners=2pt, inner sep=2pt,
              minimum width=#2, minimum height=1.6em, fill=white,
              anchor=base] (trou) {\strut\ifnumequal{1}{\UPSTIidVersionDocumentDef}%
                {\UPSTIcorrection[1]{\ttfamily #1}}{}};%
    }%
}
